% Laget automatisk av elevark/lag_elevark.py fra opplegg/statisk-elektrisitet.html. Endringer her blir overskrevet.
\documentclass[11pt]{article}
\usepackage{elevark}
\begin{document}
\elevarkhode{Fysikk 2 · Elektriske felt og krefter}{Statisk elektrisitet}{https://fysikklab.no/opplegg/statisk-elektrisitet.html}{Opplegget på nett}

Gni en ballong mot håret, og den klistrer seg til veggen. Gni en plaststav med en ullklut, og den bøyer en vannstråle. Her skal dere finne ut hvilke gjenstander som blir positive og hvilke som blir negative, hvorfor også nøytrale ting tiltrekkes, og til slutt måle hvor stor ladningen faktisk er.

\section*{Hensikt}

\begin{itemize}
\item Vise at gniing flytter elektroner fra en gjenstand til en annen, og at det finnes to typer ladning.
\item Forklare hvorfor ladde gjenstander også tiltrekker nøytrale gjenstander (influens og polarisering).
\item Bruke Coulombs lov og kraftbalanse til å anslå ladningen på to kuler.
\end{itemize}

\section*{Utstyr}

\begin{itemize}
\item Staver av glass (eller akryl) og av plast (PVC, ebonitt eller en plastlinjal), og kluter av silke og ull (eller pels).
\item Ballonger, sugerør, små papirbiter og en vannkran med tynn, jevn stråle.
\item Et elektroskop (kjøpt eller hjemmelaget av en glasskrukke, en spiker og to strimler aluminiumsfolie).
\item Til stasjon 5: to like, lette kuler (for eksempel isopor dekket med aluminiumsfolie, eller kuler av sammenklemt folie) på tynne tråder, et stativ og en linjal. En vekt som måler hundredels gram, eller kjent masse.
\item Forsøkene virker best når lufta er tørr, for eksempel om vinteren. Er det fuktig, lekker ladningen raskt bort.
\end{itemize}

\section*{Stasjon 1: Positiv eller negativ?}

\begin{enumerate}
\item Heng et sugerør i en tråd. Gni sugerøret med ullkluten, så det blir negativt ladd. Det er testgjenstanden deres.
\item Gni en plaststav med ull, og hold den nær sugerøret. Tiltrekkes eller frastøtes det?
\item Gni en glassstav med silke, og gjør det samme.
\item Prøv andre kombinasjoner: ballong og hår, plastlinjal og silke, kluten selv. Skriv i en tabell hvilken ladning hver gjenstand får.
\end{enumerate}

\section*{Stasjon 2: Nøytrale ting tiltrekkes}

\begin{enumerate}
\item Gni en ballong mot håret eller en ullgenser, og hold den mot veggen. Hvorfor blir den hengende?
\item Hold den ladde staven over noen små papirbiter. Hva skjer når bitene har vært borti staven en stund?
\item Hold den ladde staven nær en tynn vannstråle. Prøv både den positive og den negative staven. Bøyer strålen seg samme vei?
\end{enumerate}

\section*{Stasjon 3: Elektroskopet}

\begin{enumerate}
\item Før en ladd stav nær kula på toppen av elektroskopet, uten å røre. Hva skjer med bladene? Hva skjer når staven fjernes?
\item Rør kula med den ladde staven. Hva skjer nå når staven fjernes?
\item Før staven nær uten å røre, og rør kula med en finger samtidig. Fjern fingeren først og så staven. Hvilken ladning har elektroskopet nå? Test med den andre staven.
\end{enumerate}

\section*{Stasjon 4: Like og ulike kuler}

\begin{enumerate}
\item Heng de to kulene ved siden av hverandre. Lad begge med den samme staven. Hva skjer?
\item Lad den ene med plaststaven og den andre med glassstaven. Hva skjer nå?
\item Ta på den ene kula med fingeren. Hva skjer da?
\end{enumerate}

\section*{Stasjon 5: Mål ladningen}

\begin{enumerate}
\item Vei kulene, eller bruk oppgitt masse. Heng dem i like lange tråder fra samme punkt, og mål trådlengden $L$ fra festepunktet til midten av kulene.
\item Lad begge kulene med samme stav og like mye gniing, slik at de får tilnærmet like stor ladning. Rør dem gjerne mot hverandre etterpå, så ladningen fordeler seg likt.
\item Mål avstanden $d$ mellom midtpunktene til kulene, raskt, med linjalen bak kulene. Gjenta målingen hvert halve minutt mens ladningen lekker ut.
\end{enumerate}

\section*{Oppgaver}

\begin{oppgaveliste}
\oppgave \fase{gjett}{Gjett}Du gnir en plaststav med en ullklut. Får staven og kluten samme ladning, motsatt ladning, eller blir bare staven ladd?
\svar
\oppgave \fase{forklar}{Forklar}Lag en tabell over gjenstandene fra stasjon 1 og ladningen de fikk. Hva skjer egentlig når to gjenstander gnis mot hverandre? Hva er det som flytter seg?
\svar
\oppgave \fase{forklar}{Forklar}Ballongen henger på veggen, og vannstrålen bøyer seg mot staven uansett fortegn. Forklar med influens og polarisering. Tegn ladningene.
\svar
\oppgave \fase{forklar}{Forklar}Forklar det dere så med elektroskopet i stasjon 3. Hvorfor får elektroskopet motsatt ladning av staven når dere lader det ved influens?
\svar
\oppgave \fase{test}{Mål}Gjør stasjon 5, eller trykk på «Fyll inn eksempel». Hvor stor er ladningen, og hvor mange elektroner tilsvarer den? Hvor raskt lekker ladningen ut?
\svar
\oppgave \fase{test}{Regn}Tegn kreftene på én kule (tyngde, snordrag og elektrisk kraft), og vis at $F_{\mathrm{e}} = mg \tan  θ$ og $q = d \cdot \sqrt{F_{\mathrm{e}}/k}$.
\svar
\oppgave \fase{forklar}{Vurder}Hvilke forenklinger bygger beregningen på? Når blir de dårligst? Hvorfor er det vanskelig å gjøre målingen med ballonger?
\svar
\oppgave \fase{utfordring}{Utfordring}Dere går over et teppe og får et støt når dere tar i dørhåndtaket. Anta at kroppen har fått ladningen 1~µC. Hvor mange elektroner er det, og hvor stor del av elektronene i kroppen er det omtrent? (Kroppen har i størrelsesorden $10^{28}$ elektroner.)
\svar
\end{oppgaveliste}

\end{document}
